\begin{abstract}
The splitting-off operation in undirected graphs is a fundamental reduction operation that detaches all edges incident to a given vertex and adds new edges between the neighbors of that vertex while preserving their degrees. Lov\'{a}sz \cite{Lov74, Lovasz-problems-book} and Mader \cite{Mad78} showed the existence of this operation while preserving global and local connectivities respectively in graphs under certain conditions. These results have far-reaching applications in graph algorithms literature \cite{Lov76, Mad78, Frank93:Applications-of-submod-functions, FK02, Kiraly-Lau, Fra92, GB93, Fra94, JFJ95, Frank-book, nagamochi_ibaraki_2008_book, NNI97, HW96, Goe01, Jor03, Kri03, JMS03, CFLY11, BHTP03, Lau07, CS08, NZ20, BN23}. 
In this work, we introduce a splitting-off operation in hypergraphs. We show that there exists a local connectivity preserving complete splitting-off in hypergraphs and give a strongly polynomial-time algorithm to compute it in weighted hypergraphs. We illustrate the usefulness of our splitting-off operation in hypergraphs by showing two applications: (1) we give a constructive characterization of $k$-hyperedge-connected hypergraphs and (2) we give an alternate proof of an approximate min-max relation for max Steiner rooted-connected orientation of graphs and hypergraphs (due to Kir\'{a}ly and Lau \cite{Kiraly-Lau}). Our proof of the approximate min-max relation for graphs circumvents the Nash-Williams' strong orientation theorem and uses tools developed for hypergraphs.
\end{abstract}